\scanpage{50}
\deleted{\unsure{Autre façon|Autre fois}}

Ceci ne nous a pas donné de façon très explicite de classifier les \textquotedbl courbes \inserted{\unsure{simples|simple}\textquotedbl{}} de genre $0$\textquotedbl\ sur $S$, munies d'une section et d'une trisection non ramifiée, disjointe
% slanted
\margin{fonctoriel\quad\tikz[baseline,overlay]{\draw[thin] (.1em,4.2ex) .. controls (.13em,2ex) and (.36em,-1ex) .. (.12em,-5.5ex);}}
[ \deleted{\unsure{$\pi$|$X$}} = revêtement \deleted{\unsure{$\mathfrak S_3$}} \unsure{non|don} \deleted{\unsure{g|G}} ramifié \inserted{$S'$} de groupe $\mathfrak S_3$ de $S$, \unsure{avec|par|pour} une section \deleted{\unsure{de}} de
$\deleted{\unsure{X=S'\times_{(S,\mathfrak S_3)}P^1_S}}S'\times_{(S,\mathfrak S_3)}P^1_S$ \inserted{$t$ disjoint de $0,1,\infty$} invariante par le groupe $\mathfrak S_3$, i.e. \deleted{une section}
\par\noindent\begin{minipage}[t]{.21\linewidth}\vspace{0pt}\scriptsize
\diagram{\begin{tikzcd}[ampersand replacement=\&,column sep=1.5em,row sep=1em]
X\arrow[d]\&P^1_{S'}\arrow[l]\arrow[d]\\
S\&S'\arrow[l,"\mathfrak S_3"]
\end{tikzcd}}
\figure{\begin{tikzpicture}[x=1em,y=1ex]\draw[thin] (0,2.2) .. controls (-.1,.8) and (.1,-1.7) .. (.25,-2.45) .. controls (.3,-2.8) and (.58,-3.45) .. (.69,-3.8);\end{tikzpicture}}
\end{minipage}\hfill\begin{minipage}[t]{.74\linewidth}\vspace{0pt}
$t\in\Gamma(S',\underline{\mathcal O}_{S'})$, inversible et telle que $t-1$ soit inversible, et satisfaisant à
\[
\sigma.t=\deleted{\unsure{\varphi}}\bar\sigma(t)\qquad \sigma\in\mathfrak S_3
\]
i.e.
\[
\left\{\begin{aligned}
\lambda t&=\frac1t\\
\mu t&=1-t
\end{aligned}\right.
\]


\end{minipage}

\uline{Exemple} \quad $S=\operatorname{Spec}(O)$, $O$ anneau local \uline{complet} à corps résiduel alg\textsuperscript{t} clos \deleted{\unsure{(de car. $\ne2,3$)}} \inserted{(de car. $\unsure{\ne2,3|\ne1,3}$)}.
% slanted
\margin{\uuline{non} \uline{fonctoriel}}
Alors \deleted{\unsure{$X$|$X'$} \illegible{1} $\unsure{S}$ \illegible{2}} $\pi_1(S)=0$, donc la situation équivaut, \inserted{\unsure{isom. non canon. près|isom. non canonique près|Dém. non canon. près},} à la donnée d'un $t\in O^*$ tel que $\deleted{\unsure{t(t-1)|t(t-2)}\unsure{\in O^*|\in}}t(t-1)\in O^{\unsure{+|*}}$, étant entendu que $t,t'$ donnent \unsure{des|deux} syst. isomorphes \deleted{si et} s. s. s. $t'$ est de la forme $\bar\sigma(t)\quad\sigma\in\mathfrak S_3$.
% archive sheet number: 49
